%%=============================================================================
%% Methodologie
%%=============================================================================

\chapter{\IfLanguageName{dutch}{Methodologie}{Methodology}}%
\label{ch:methodologie}
Om de onderzoeksvraag te beantwoorden werden vijf deelvragen beantwoord: (a) Waarom is er nood aan modernisatie van mainframes?; (b) Wat doet automatisatie op mainframes?; (c) Wat zijn de voor- en nadelen van Rexx?; (d) Wat zijn de voor- en nadelen van Python op z/OS?; (e) In hoeverre kan Python nu al Rexx vervangen? De vijf deelvragen werden benaderd via twee onderzoekstechnieken: een literatuurstudie en een vergelijkende studie. De onderzoekstechnieken bestaan uit verschillende stappen die verdeeld werden in drie fases. De eerste fase was de literatuurstudie, die een antwoord bood op de eerste twee deelvragen, en de basis legde voor deelvragen 3 en 4. De tweede fase was een requirementsanalyse. Die hielp met (1) het identificeren van veelgebruikte functionaliteiten en use cases waarmee Rexx en Python kunnen worden vergeleken; (2) criteria waaraan beide talen moeten voldoen. Tot slot, in de derde fase, volgden de PoC's van elke use case, die samen meer inzicht gaven in de voor- en nadelen van beide talen (deelvraag 3 en 4), en die een antwoord boden op de vraag of Python nu al Rexx kan vervangen (deelvraag 5).

%% TODO: In dit hoofstuk geef je een korte toelichting over hoe je te werk bent
%% gegaan. Verdeel je onderzoek in grote fasen, en licht in elke fase toe wat
%% de doelstelling was, welke deliverables daar uit gekomen zijn, en welke
%% onderzoeksmethoden je daarbij toegepast hebt. Verantwoord waarom je
%% op deze manier te werk gegaan bent.
%% 
%% Voorbeelden van zulke fasen zijn: literatuurstudie, opstellen van een
%% requirements-analyse, opstellen long-list (bij vergelijkende studie),
%% selectie van geschikte tools (bij vergelijkende studie, "short-list"),
%% opzetten testopstelling/PoC, uitvoeren testen en verzamelen
%% van resultaten, analyse van resultaten, ...
%%
%% !!!!! LET OP !!!!!
%%
%% Het is uitdrukkelijk NIET de bedoeling dat je het grootste deel van de corpus
%% van je bachelorproef in dit hoofstuk verwerkt! Dit hoofdstuk is eerder een
%% kort overzicht van je plan van aanpak.
%%
%% Maak voor elke fase (behalve het literatuuronderzoek) een NIEUW HOOFDSTUK aan
%% en geef het een gepaste titel.

\section{Literatuurstudie}
In de literatuurstudie werd getracht om een overzicht te geven van de huidige situatie in het veld van mainframe. Daarvoor werd er eerst naar het verleden gekeken en hoe Rexx zo een grote scripting taal van mainframe is geworden. Daarna werd er gekeken naar de opkomst van Python en zijn populariteit. Uiteindelijk werd een analyse gemaakt van bestaand onderzoek dat de twee talen vergelijkt. Deze vergelijking, waaruit bleek dat beide talen hun voor- en nadelen hebben, vormde de basis voor het verdere onderzoek in deze bachelorproef. In de literatuurstudie werden de eerste twee deelvragen (Waarom is er nood aan modernisatie op mainframes? \& Wat doet automatisatie op mainframes?) volledig beantwoord. Daarnaast werden de derde en vierde deelvraag (Wat zijn de voor- en nadelen van Rexx? \& Wat zijn de voor-en nadelen van Python op z/OS?) ook al deels beantwoord. Die laatste vergelijking in de literatuur werd echter voornamelijk gebruikt als ondersteuning bij het bepalen van de requirements die nodig waren voor een succesvolle proof-of-concept. De resultaten van deze fase zijn terug te vinden in hoofdstuk \ref{ch:stand-van-zaken}.

\section{Requirementsanalyse}
Een requirementsanalyse werd uitgevoerd om te bepalen welke veelgebruikte functies en use cases werden gebruikt om de vergelijking tussen Rexx en Python te maken. Daarnaast werden in deze fase de succescriteria bepaald waaraan beide talen werden getoetst om tot een conclusie te komen over de beste optie voor elke use case. Om de best mogelijke use cases en succescriteria te bepalen, werden verschillende interviews afgenomen met ervaren systeemprogrammeurs, en werd een vragenlijst opgesteld om een brede input te krijgen van mensen in het werkveld. Nadien werd er met behulp van de MoSCoW-methode een duidelijk beeld geschetst van de vereisten waaraan beide talen moeten voldoen om te slagen in een use case. Zo werd de voorkeur voor één van de twee talen in die specifieke use case bepaald. 

\subsection{Interviews}
De interviews werden afgenomen met ervaren systeemprogrammeurs. De informatie uit de interviews werd gebruikt om een goed beeld te krijgen van hoe mainframers de modernisatie op het mainframe ervaren, alsook om input te krijgen over de use cases en de requirements die werden gekozen voor de PoC.

De uitgeschreven interviews zijn opgenomen in de bijlage \ref{ch:Interviews}.

\subsection{Enquête}
De enquête werd opgesteld met het doel om input te krijgen van zoveel mogelijk systeemprogrammeurs, en dit met een lage barrière. Vooral de succescriteria werden bevraagd in de enquête, aangezien daar verschillende meningen over kunnen zijn en dit onderzoek een objectieve analyse wil uitvoeren. De enquête werd gedeeld in de Discord community van 'System Z Enthusiasts' en op Linkedin. De enquête is te vinden in de bijlage \ref{ch:Enquête}.

\section{Proof-of-Concept}
In de Proof-of-Concept werd de vergelijkende studie uitgevoerd. Hierbij werden verschillende kleinere PoC's opgesteld aan de hand van de eerder bepaalde use cases en requirements. Er werden verschillende use cases uitgewerkt, om zo een breed beeld te krijgen van de verschillen tussen beide talen, met elk hun sterktes en zwaktes. Zo konden er use cases worden uitgewerkt waar de ene taal een duidelijk voordeel heeft ten opzicht van de andere. Bij elke PoC werd beargumenteerd waarom dit een belangrijke use case is, en wat de succescriteria zijn waaraan een oplossing moet voldoen. Daarnaast werd bij elke PoC de gebruikte code getoond en uitgelegd. En ook de bevindingen werden toegelicht aan de hand van de vooropgestelde succescriteria.

\section{Resultaten}
In de resultaten werd er gereflecteerd over de scope van het onderzoek en enkele beperkingen die mogelijks invloed hebben gehad op de resulaten van dit onderzoek. Daarnaast werden de bevindingen uit de verschillende PoC's samengenomen en geanalyseerd en werd gereflecteerd over welke use cases in toekomstig onderzoek zouden kunnen worden uitgevoerd.

\section{Conclusie}
In de conclusie werd teruggekoppeld naar de onderzoeksvragen vanuit de bevindingen en resultaten van het onderzoek. Ook werd een reflectie toegevoegd over de resultaten met het oog op het ruimere ecosysteem van mainframe en de toekomst ervan. Op basis daarvan werd een advies geformuleerd voor bedrijven over de aanpak van automatisatie.